\documentclass[10pt,a4paper]{amsart}
\usepackage[margin=19mm]{geometry}
\usepackage{amsmath,amssymb,mathtools}
\usepackage[scr=rsfs,cal=euler]{mathalfa}
\usepackage{xfrac}
\usepackage[unicode,hidelinks,bookmarks=false]{hyperref}
\newcommand{\ie}{i.e.\ }
\newcommand{\cf}{cf.\ }
\newcommand{\resp}{respectively\ }
\newcommand{\Subsection}{\subsection}
\providecommand{\nobreakdash}{\nobreak\hbox{-}}
\setlength{\emergencystretch}{2em}
\multlinegap0pt
\usepackage{algorithm}\usepackage{algpseudocode}
\usepackage{amssymb}
\usepackage{mathtools}
\usepackage{bm}
\usepackage{xcolor}
\usepackage[shortlabels]{enumitem}
\newtheorem{theorem}{Theorem}
\newcommand{\Pcal}{\mathcal P}
\newcommand{\e}{\mathrm e}
\title{Least Variability in a Polynomial-Square Class\\ of Rational Kernels}
\author{Maria Laura Battagliola, Oscar Peralta}
\date{Source: arXiv:2608.29143v1 (CC BY 4.0)}
\begin{document}
\maketitle

\noindent\emph{Source and licence.} Least Variability in a Polynomial-Square Class\\ of Rational Kernels. Version arXiv:2608.29143v1 (CC BY 4.0), distributed by arXiv under the Creative Commons Attribution 4.0 International licence, http://creativecommons.org/licenses/by/4.0/, as recorded in the arXiv licence metadata for this exact version and shown on the abstract page https://arxiv.org/abs/2608.29143v1. Full licence notice: \texttt{licenses/arxiv-2608.29143-CC-BY-4.0.txt}.\par\smallskip

\noindent\textit{Editorial scope.} Counted unit: the article's theorem var:thm:first-variation (source0.tex lines 784--801). The article's complete original proof of it is delivered with it (source0.tex lines 803--839, one \texttt{proof} environment of 37 lines). No mathematics is written, repaired or re-derived: the statement and the proof are the article's own lines, in the article's own notation, and no retained line is substituted or re-flowed.  Retained uncounted as the source-local context the proof needs: lines 735--744.  Nothing was omitted from any retained span, so the package carries no omission record.  No retained line contains a citation, so the wrapper carries no bibliography and no bibliography entry.  No retained line contains a figure environment, an image-inclusion command or drawing code, so no image is delivered and the package declares no asset dependency.  Editorial formatting only, in addition: an AMS \texttt{amsart} preamble that restates the article's own declarations and macros used by the retained lines, namely the environments \texttt{theorem} on separate counters, together with the standard packages those lines need.  The retained environments print in this document as Theorem 1 in order of appearance, whereas the article prints the same items with the numbering of its own full document.  The complete native source of the article as distributed in the arXiv e-print for this exact version is bundled unchanged as \texttt{sources/208830/source0.tex}.
\begin{equation}\label{var:eq:J-def}
\mathcal J(a,\mu) =
\frac{
\displaystyle\int_0^\infty (t-\mu)^2a(t)^2e^{-t},dt
}{
\displaystyle\mu^2\int_0^\infty a(t)^2e^{-t},dt
},
\qquad
0\ne a\in\Pcal_m,\quad \mu>0.
\end{equation}

\begin{theorem}[First variation in polynomial directions]\label{var:thm:first-variation}
Let \((a_*,\mu_*)\) be a joint minimizer of \(\mathcal{J}\) in \eqref{var:eq:J-def}, and let \(\rho_*=\mathcal{J}(a_*,\mu_*)\). Define
\begin{equation}\label{var:eq:eta-def}
    \eta_*
    :=\rho_*\mu_*^2
    =\min_{0\ne a\in\Pcal_m}\Lambda_{\mu_*}(a),
\end{equation}
with $a_*$ realizing this minimum, and set
\begin{equation}\label{var:eq:R-def}
    R_*(t)=(t-\mu_*)^2-\eta_*.
\end{equation}
Then
\begin{equation}\label{var:eq:Euler-integral}
    \int_0^\infty R_*(t)a_*(t)b(t)\e^{-t}\,dt=0
    \qquad
    \text{for every }b\in\Pcal_m.
\end{equation}
\end{theorem}

\begin{proof}
For \(b\in\Pcal_m\), consider the path \(a_\varepsilon=a_*+\varepsilon b\), admissible since \(a_*\ne0\) keeps \(a_\varepsilon\) nonzero for small \(\varepsilon\). Holding \(\mu=\mu_*\) fixed and writing
\begin{align*}
    N(\varepsilon)
    &=\int_0^\infty
      (t-\mu_*)^2a_\varepsilon(t)^2\e^{-t}\,dt,
      &
    D(\varepsilon)
    &=\mu_*^2\int_0^\infty
      a_\varepsilon(t)^2\e^{-t}\,dt,
\end{align*}
we have \(\mathcal{J}(a_\varepsilon,\mu_*)=N(\varepsilon)/D(\varepsilon)\). Since \((a_*,\mu_*)\) is a joint minimizer, this ratio has a local minimum at \(\varepsilon=0\), so
\begin{equation}\label{var:eq:stationarity-epsilon}
    N'(0)D(0)-N(0)D'(0)=0.
\end{equation}
Expanding \(a_\varepsilon^2=a_*^2+2\varepsilon a_*b+\varepsilon^2b^2\) gives
\begin{align}
    N'(0)
    &=2\int_0^\infty
      (t-\mu_*)^2a_*(t)b(t)\e^{-t}\,dt,
      \label{var:eq:N-prime}\\
    D'(0)
    &=2\mu_*^2\int_0^\infty
      a_*(t)b(t)\e^{-t}\,dt.
      \label{var:eq:D-prime}
\end{align}
Since \(\rho_*=N(0)/D(0)\), \eqref{var:eq:stationarity-epsilon} reduces to \(N'(0)-\rho_*D'(0)=0\). Inserting \eqref{var:eq:N-prime}--\eqref{var:eq:D-prime}, dividing by two, and combining the integrals gives
\begin{align*}
    0
    &=\int_0^\infty
      \left[(t-\mu_*)^2-\rho_*\mu_*^2\right]
      a_*(t)b(t)\e^{-t}\,dt
    =\int_0^\infty
      R_*(t)a_*(t)b(t)\e^{-t}\,dt,
\end{align*}
which proves \eqref{var:eq:Euler-integral}.
\end{proof}

\end{document}
